\section{Rational dilations of arithmetic frames and an exact completion obstruction}
\label{sec:rational-dilations}

Concatenation produces the partial matrices used in the arithmetic
branch. It is useful to understand the remaining completion problem
as an explicit orthogonal dilation. The elementary calculation below
separates real or rational completion from the sign constraints.
No novelty is claimed for the linear algebra or the quaternion
sum-of-four-squares construction.

Let $a<b$ be positive multiples of four, let $A$ be an order-$a$
Hadamard matrix, and partition an order-$b$ Hadamard matrix as
$H_b=\binom{B}{E}$, with $B$ having $a$ rows. Put $d=b-a$ and $n=a+b$.
Then $P=[A\ B]$ is an $a$-row sign frame of width $n$.

\begin{theorem}[An exact residual parameterization]
\label{dil:parameterization}
The real matrices
\[
 Z=\begin{pmatrix}A&B\\B^T&D\end{pmatrix},\qquad ZZ^T=nI_n,
\]
are exactly those with
\[
 \boxed{D=\frac1b\left(-B^TA^TB+E^TVE\right),\qquad
                         VV^T=nI_d.}
\]
In particular a rational such $Z$ always exists. Its only possibly
nonsign entries are in the lower-right $b$-square block.
A sign-valued completion with this prescribed transpose corner
exists if and only if some such $V$ makes every entry of $D$ equal
to $+1$ or $-1$.
\end{theorem}
\begin{proof}
The top block already has Gram $nI_a$. The remaining equations are
\[
 AB+BD^T=0,\qquad B^TB+DD^T=nI_b.
\]
Conjugate $D$ by the orthogonal matrix
$W=\binom{B}{E}/\sqrt b$. The first equation forces
\[
 WDW^T=\begin{pmatrix}-A^T&F\\0&V\end{pmatrix}.
\]
The second requires its product with its transpose to be
$\operatorname{diag}(aI_a,nI_d)$.
Since $A^TA=aI_a$, this forces $FF^T=0$, hence $F=0$,
and leaves exactly $VV^T=nI_d$. Transforming back gives the formula.

For rational existence write $n=u_0^2+u_1^2+u_2^2+u_3^2$ by
Lagrange's four-square theorem. The quaternion multiplication matrix
\[
 Q=\begin{pmatrix}
 u_0&u_1&u_2&u_3\\
 -u_1&u_0&-u_3&u_2\\
 -u_2&u_3&u_0&-u_1\\
 -u_3&-u_2&u_1&u_0
 \end{pmatrix}
\]
satisfies $QQ^T=nI_4$. Since $d$ is a multiple of four, take $V$
block diagonal with $d/4$ copies of $Q$. It is integral, and the
displayed formula gives a rational $D$. The sign-valued criterion
is then simply the definition of a Hadamard matrix with these
prescribed three blocks.
\end{proof}

The free real parameter lies in a scaled $O(d)$, rather than $O(n)$.
This is a reduction of the completion variables, not a proof of
flatness. In particular rational scaled orthogonality supplies no
general sign-rounding theorem.

\subsection{A structured order-twenty obstruction and its smallest repair}
Take $A=H_8$ in the standard Sylvester order, and let $B$ be the
first eight rows of the normalized Paley type-II $H_{12}$ constructed
from the prime five. The exact frame $[A\ B]$ has eight orthogonal
rows of length twenty. It cannot be completed to a Hadamard matrix
while preserving all eight rows.

Here is a finite exhaustive certificate of that assertion.
Its row span modulo five has dimension eight. The four-dimensional
hyperbolic quotient therefore has exactly twelve self-dual
extensions, the two rulings of the preceding pass's $6$-square grid.
The checker reconstructs all twelve codes and exhausts each
systematic sign assignment. Each code has exactly twelve sign
classes, and those classes are real orthogonal because the
prescribed eight rows fix their parity. Thus the exact maximum
orthogonal frame preserving these eight rows is twelve, not twenty.
The row residues of any full Hadamard completion would lie in one
of these self-dual extensions, contradicting its sign count.

Deleting one prescribed row does permit a full completion.
The checker enumerates modular candidates orthogonal to the retained
seven rows, makes an explicit self-dual neighbor pivot, and obtains
twenty pairwise orthogonal sign classes preserving exactly seven of
the eight originals. The certificate and row index are in
\path{dilations/checks.json}. The minimum original-row loss for a
full completion is consequently exactly one.
This is a known-order example; it proves neither a uniform bounded
repair theorem nor failure of the arithmetic construction at other orders.

The continuous family is nonempty despite the sign obstruction.
Using the quaternion parameters $(4,2,0,0)$ gives $VV^T=20I_4$.
The checker verifies the full rational $20$-square Gram identity,
with lower-right entries ranging from $-2$ to $8/3$ and
\[
                    \sum_{ij}D_{ij}^4-144=\frac{6976}{27}.
\]
Rounding its entries to signs does not produce a Hadamard matrix.

\subsection{A quantitative flatness gap when the prescribed rows cannot complete}
\begin{proposition}[A rounding threshold for real completions]
\label{dil:flatness-gap}
Let $n$ be a positive multiple of four and let $Z$ be any real matrix
with $ZZ^T=nI_n$. Define
\[
 \Delta=\sum_{ij}(Z_{ij}^2-1)^2=\sum_{ij}Z_{ij}^4-n^2.
\]
If
\[
 \Delta<(\sqrt{n+2}-\sqrt n)^2,
\]
rounding each entry to a sign gives a Hadamard matrix.
Consequently if some prescribed sign rows cannot be completed,
every real completion preserving them has
$\Delta\ge(\sqrt{n+2}-\sqrt n)^2>1/(n+2)$.
\end{proposition}
\begin{proof}
Let $H$ be the entrywise sign rounding (choose $+1$ at zero)
and write $F=H-Z$. The scalar inequality
$(|z|-1)^2\le(z^2-1)^2$ gives $\|F\|_F\le\sqrt\Delta$.
Since $\|Z\|_{\rm op}=\sqrt n$,
\[
 \|HH^T-nI_n\|_{\rm op}
 \le2\sqrt n\,\sqrt\Delta+\Delta<2.
\]
The diagonal is zero, and every off-diagonal entry of this error
is an even integer, since $n$ is even and $H$ is a sign matrix.
The operator norm bounds every entry in magnitude, so all entries
vanish. Prescribed sign rows remain unchanged by rounding.
Finally
$(\sqrt{n+2}-\sqrt n)^2
=4/(\sqrt{n+2}+\sqrt n)^2>1/(n+2)$.
\end{proof}

For the order-twenty frame this gives an explicit positive quartic
gap for every continuous completion preserving the eight rows.
A continuous optimizer cannot remove that gap without changing
prescribed rows. This pinpoints a real limitation of feeding an
arithmetic partial frame into a monotone completion method.
